\section{Decision-evaluation protocol}
\label{sec:protocol}

\subsection{Graph-versus-MLP decision}

The unit of evaluation is a dataset--seed pair $u$. Within each unit, validation performance selects one feature-only MLP and one graph model from GCN, GAT, GraphSAGE, H2GCN, LINKX, and GPR-GNN. Every architecture receives the same four-trial tuning budget, and exact validation ties are broken by model identifier. Let $a_u^{\mathrm{M}}$ and $a_u^{\mathrm{G}}$ denote the held-out test accuracies of the selected MLP and graph model. The evaluator reads each selected test result once; no diagnostic can affect training or hyperparameter selection. The resulting actions are portfolios: graph is the validation-selected member of six architectures, whereas MLP is the validation-selected configuration of one architecture.

We use a frozen practical margin $\delta=0.01$ to define the target action:
\begin{equation}
y_u =
\begin{cases}
\mathrm{G}, & a_u^{\mathrm{G}}-a_u^{\mathrm{M}}>\delta,\\
\mathrm{M}, & a_u^{\mathrm{G}}-a_u^{\mathrm{M}}\leq\delta.
\end{cases}
\label{eq:decision_target}
\end{equation}
Every rule uses this asymmetric target: graph must exceed MLP by more than one percentage point, and ties favor MLP.

\subsection{Per-architecture training and selection}

Table~\ref{tab:frozen_training} records the training grid. Each PyTorch trial~\citep{paszke2019pytorch} minimizes cross-entropy with Adam and retains the checkpoint with highest validation accuracy, then lowest validation loss and earliest epoch. Training stops after 100 stale epochs or 500 epochs in total. The grid crosses learning rates $\{0.01,0.005\}$ and dropout rates $\{0.5,0.7\}$, with width 64 and weight decay $5\times10^{-4}$. Each architecture independently selects one of its four trials and evaluates the chosen checkpoint once on the held-out test set. The six-architecture graph portfolio therefore searches 24 trial configurations, compared with four for MLP.



The split targets 60/20/20 for train/validation/test and is regenerated deterministically for seeds 0--9. Each class must contain at least three nodes, with at least one assigned to each of the training, validation, and test sets. Integer per-class floors make the realized proportions approximate, especially for small WebKB graphs; the actual counts are reported in the reproducibility summary. Every unit stores the source commit, canonical configuration digest, processed-data checksums, split identifier, four trial records, selected checkpoint, and environment provenance. The v1 and v2 JSON specifications differ only in the pre-outcome eligibility amendment described in Section~\ref{sec:prospective}; the later Squirrel amendment changes only the execution timeout.

The frozen protocol in this section governs the original benchmark. Subsequent portfolio analyses are post-hoc, and the additional preprocessing and training diagnostics use separately recorded configurations and information boundaries. Their status and scope are identified where they are reported.

\subsection{Decision-time information and policies}

Full node features and graph structure are available because they are model inputs. Label-dependent statistics use the subgraph induced by training nodes. Edge homophily $h_1$ is the label-agreement fraction on its edges. For the two-hop statistic $h_2$, a center with training-subgraph degree $d_T$ is sampled with probability proportional to $d_T(d_T-1)$, then two distinct neighbors are sampled uniformly without replacement. All three nodes therefore belong to the training set. The estimate is the mean endpoint-label agreement over 50,000 sampled ordered length-two paths, using the split seed plus 200,000 as the NumPy generator seed. Paths may repeat across draws, and adjacent endpoints are included when they form a length-two path; this is a path-weighted statistic, not strict distance-two homophily. If no eligible edge or path exists, the corresponding statistic is missing. Test labels, test accuracy, and statistics computed from test labels are forbidden diagnostic inputs; the evaluator rejects records that violate this boundary.

This is an offline audit of predictive decision utility: all candidate models are trained to supply common outcomes, and regret measures the test accuracy lost by selecting among them. Table~\ref{tab:diagnostic_policies} distinguishes each policy's required information and decision stage. An early family choice still requires fitting and selecting a model within that family. Validation selection provides a reference for decisions made after both families have been tuned. The evaluation measures accuracy loss and coverage; training-cost savings and deployment efficiency are outside its measured outcomes.

Table~\ref{tab:diagnostic_policies} lists the nine policies. The Combined rule has an exploratory origin: its exact $0.55/0.40$ thresholds were fixed in the specification and scorer on 8 August 2026, after inspection of legacy aggregates and before generation of the primary records. The final 110 units were not used to adjust them. Appendix~\ref{app:execution_details} retains the threshold provenance; the scope is a pre-specified primary analysis rather than whole-project preregistration. The policy chooses graph when training-only edge homophily is at least $0.55$; below that threshold, it chooses MLP when MLP validation accuracy is at least $0.40$ and otherwise abstains. Missing inputs also cause abstention. For full-set evaluation, every abstention uses the same predeclared MLP fallback, while coverage remains separately visible.

All inputs are complete in these records, so Combined's full-set action simplifies to graph if $h_1\geq0.55$ and MLP otherwise. The $0.40$ MLP-validation cutoff changes explicit-action labels, abstention coverage, and the confidence curve. Under MLP fallback, it leaves full-set accuracy and regret unchanged. We retain the name Combined for the historical rule before fallback.

The exact nine-policy definitions and curve-ordering rules are retained in Appendix~\ref{app:frozen_policy_details}.

\subsection{Outcomes and inference}

Selection accuracy measures agreement with Eq.~\eqref{eq:decision_target} after the common fallback, and coverage is the fraction of non-abstaining decisions. Because an incorrect action can have negligible or substantial cost, the primary loss is oracle-portfolio regret,
\begin{equation}
r_u(\widehat{y})=\max\{a_u^{\mathrm{G}},a_u^{\mathrm{M}}\}-a_u^{\widehat{y}},
\label{eq:oracle_regret}
\end{equation}
where $a_u^{\widehat{y}}$ is the test accuracy of the chosen portfolio after fallback. Selection accuracy uses the one-point practical-margin target in Eq.~\eqref{eq:decision_target}; regret measures raw accuracy lost relative to the better selected portfolio. Choosing MLP for a sub-margin graph gain therefore agrees with the target but incurs a small positive regret. We report mean full-set regret for all units and covered-set regret for non-abstaining units. Descriptive regret--coverage curves rank eligible decisions by each policy's frozen confidence score and plot covered-set raw-accuracy regret. They leave the frozen target and primary full-set analysis unchanged.

Seeds remain paired within each dataset. We average regret over seeds, then use dataset means for percentile-bootstrap intervals and two-sided sign-flip tests. The sign-flip reference distribution is exact when dataset-level differences are jointly invariant to the enumerated sign changes. Independent differences symmetric about zero suffice; a zero mean or pairing alone does not. Holm correction controls family-wise error across the eight predeclared comparisons with Combined when the constituent p-values are valid. The convenience-selected datasets may be dependent and unrepresentative, so both intervals and tests have a conditional, assumption-dependent interpretation. Non-significance does not establish equivalence.

The frozen paper-level rule permits an incremental-value claim only if Combined reduces full-set regret relative to simple baselines without relying on lower coverage. It also requires the uncertainty interval to exclude zero in the favorable direction for the predeclared comparison. Each of the other eight policies in Table~\ref{tab:diagnostic_policies} is compared with Combined using an unadjusted 95\% dataset-level percentile-bootstrap interval. The specification does not single out always-graph or resolve whether its criterion must hold for all baselines. We use always-graph as the focal comparator in this report. Positive baseline-minus-Combined differences favor Combined. Full-set scoring includes every abstention through the declared fallback; uncovered units cannot be omitted to meet the criterion. We report uncertainty intervals and Holm-adjusted tests separately, since significance alone does not satisfy the regret and coverage requirements.

The prospective specification, immutable per-unit schema, analysis code, bootstrap seed, permutation seed, and paper-level decision rule were fixed before assembling or evaluating the new comparative outcomes. Section~\ref{sec:prospective} reports two outcome-independent protocol amendments and keeps the earlier feasibility artifacts outside the final analysis.

\subsection{Diagnostic utility at the action level}
\label{sec:action_audit}

To compare diagnostics, we first resolve abstentions into effective actions. Let $q_u\in\{\mathrm{G},\mathrm{M}\}$ be that action and let $\mathbb{E}_D$ average seeds within each dataset and then weight datasets equally. Write $g_u=a_u^{\mathrm{M}}-a_u^{\mathrm{G}}$ and $x_+=\max(x,0)$. The improvement over always-graph satisfies the accounting identity
\begin{align}
R_{\mathrm{G}}-R_q
 &= \mathbb{E}_D[\mathbf{1}\{q_u=\mathrm{M}\}g_u] \\
 &= \underbrace{\mathbb{E}_D[\mathbf{1}\{q_u=\mathrm{M}\}(g_u)_+]}_{B_q:\ \text{benefit}}
 -\underbrace{\mathbb{E}_D[\mathbf{1}\{q_u=\mathrm{M}\}(-g_u)_+]}_{H_q:\ \text{harm}}.
\label{eq:benefit_harm}
\end{align}
Subtracting the regrets in Eq.~\eqref{eq:oracle_regret} gives this identity. Only MLP choices contribute: correct switches supply benefit, and switches away from a more accurate graph model supply harm. The decomposition attributes recorded loss to the policy actions.

The available improvement over always-graph is bounded by its own regret,
\begin{equation}
R_{\mathrm{G}}-R_q\leq R_{\mathrm{G}}=\mathbb{E}_D[(g_u)_+].
\label{eq:headroom_bound}
\end{equation}
Headroom bounds the absolute gain available on the recorded candidate outcomes. Dividing by positive headroom expresses a gain as a fraction of that ceiling. Absolute benefit and statistical uncertainty are reported separately. The ceiling depends on the portfolio and training outcomes.

Two rules with identical effective actions have identical full-set regret even if their scores, confidence rankings, or abstention coverage differ. Combined supplies a concrete example: under MLP fallback, changing its MLP-validation cutoff changes the explicit-action/abstention labels without changing its effective homophily-threshold action. Conversely, changing fallback is consequential only on abstaining units. These observations motivate reporting scores, explicit actions, fallback-resolved actions, and outcome losses as separate parts of the decision record.

\subsection{Reconstruction from saved outcomes}
\label{sec:replay_scope}

The retained records allow alternative choices among already evaluated predictors. Restricting a portfolio reapplies validation selection to retained architecture winners; changing a threshold or fallback selects between their saved outcomes. Training, candidate availability, and within-architecture selection stay fixed. Changes to preprocessing, training, hyperparameter generation, or upstream selection require new records whenever they change the available predictors. The 24-trial MLP extension therefore adds 20 trials per unit and evaluates the selected checkpoint.

For threshold calibration, selected-model test outcomes from the other ten datasets are meta-training targets. The held-out dataset contributes neither scores nor outcomes to cutoff selection. Reusing the collection across folds provides an internal cross-dataset evaluation; prospective transfer to an independent graph collection remains untested. The later matched-calibration analysis fixed its reporting scope in code before execution but is post-hoc relative to the inspected benchmark.


\subsection{Research tools}
\label{sec:research_tools}

Two generative AI coding assistants were used under the author's direction. OpenAI Codex sessions are recorded at least from 7 September to 9 October 2026. Assistance covered code and test development, experiment orchestration, result checks, analysis, and manuscript editing. This work included the frozen benchmark, sensitivity analyses, input-parameterization retraining, classifier-probe extension, and independent reader implementation. Anthropic Claude Code sessions are recorded from 24 September to 4 October 2026. Assistance covered implementation review, expanded MLP-search orchestration, audit scripting, integration of the audited extension evidence, and manuscript revision. The author set the research objective and authorized the computations; retained protocols and amendments identify the executed settings. Formal runs followed frozen configurations and validation-selection rules. Experimental and analysis programs computed every reported number, with outputs retained in the research artifacts. The author is responsible for verifying the AI-assisted code, analyses, and text before submission and for the study's interpretation. The AI systems are not authors.
